\Needspace{14\baselineskip}

\CNsection{5.1}{Why LTE Was Developed}

\CNchained\CNsubsection{}{Limitations of 3G (UMTS/\allowbreak{}HSPA)}

The 3rd Generation Partnership Project (3GPP) initiated the Long Term Evolution (LTE) project in 2004 to address fundamental limitations in 3G networks:

\textbf{Architectural Complexity:}

\begin{itemize}
\tightlist
\item
  3G UTRAN required a hierarchical structure: UE \ensuremath{\to} NodeB \ensuremath{\to} RNC \ensuremath{\to} CN
\item
  The Radio Network Controller (RNC) was a single point of failure and bottleneck
\item
  Complex protocol stacks between NodeB and RNC (Iub interface) added processing overhead
\item
  Soft handover required coordination between multiple NodeBs via the RNC
\end{itemize}

\textbf{Latency Issues:}

\begin{itemize}
\tightlist
\item
  3G round-trip times (RTT): 50–100~ms typical, 150+ ms under load
\item
  The RNC added processing hops in both control and user planes
\item
  Circuit-switched voice added further complexity and delay
\item
  Real-time applications (gaming, video calls) suffered noticeably
\end{itemize}

\textbf{Spectrum Efficiency:}

\begin{itemize}
\tightlist
\item
  WCDMA/\allowbreak{}HSPA approaching theoretical limits
\item
  Demand for mobile data growing exponentially (smartphones era beginning)
\item
  Need for flexible bandwidth: 1.4~MHz to 20~MHz channel widths
\end{itemize}

\textbf{All-IP Target:}

\begin{itemize}
\tightlist
\item
  3G still maintained circuit-switched domain for voice (MSC, MGW)
\item
  Dual-stack (CS + PS) increased OPEX and complexity
\item
  Industry consensus: everything must converge to IP
\end{itemize}

\textbf{Performance Targets for LTE (3GPP Release 8):}\CNbr{}\textbar{} Parameter \textbar{} Target \textbar{}\CNbr{}\textbar{}———–\textbar{}——–\textbar{}\CNbr{}\textbar{} Peak DL rate \textbar{} 100~Mbps (20~MHz) \textbar{}\CNbr{}\textbar{} Peak UL rate \textbar{} 50~Mbps (20~MHz) \textbar{}\CNbr{}\textbar{} User plane latency \textbar{} \textless{} 5~ms \textbar{}\CNbr{}\textbar{} Control plane latency \textbar{} \textless{} 100~ms (idle\ensuremath{\to}active) \textbar{}\CNbr{}\textbar{} Spectrum efficiency \textbar{} 3–4x improvement over HSPA \textbar{}\CNbr{}\textbar{} Mobility \textbar{} Up to 350 km/\allowbreak{}h \textbar{}

\CNsection{5.2}{LTE Design Principles}

LTE was designed from scratch with these guiding principles:

\CNsubsection{}{Flat Architecture}

\begin{itemize}
\tightlist
\item
  \textbf{Eliminate the RNC} — distribute its functions to the eNodeB and core network
\item
  Fewer network nodes = fewer hops = lower latency
\item
  Each eNodeB connects directly to the core network (EPC)
\end{itemize}

\CNsubsection{}{All-Packet (No Circuit Switching)}

\begin{itemize}
\tightlist
\item
  \textbf{Pure packet-switched network} — no circuit-switched domain
\item
  Voice handled via VoLTE (Voice over LTE) using IMS
\item
  Single protocol stack simplifies operations
\end{itemize}

\CNsubsection{}{No Soft Handover}

\begin{itemize}
\tightlist
\item
  3G soft handover required simultaneous connections to multiple NodeBs
\item
  LTE uses \textbf{hard handover} only — simpler, more spectrum-efficient
\item
  Handover is fast enough (\textasciitilde{}20–50~ms interruption) to be imperceptible
\end{itemize}

\CNsubsection{}{Simplified Protocol Stack}

\begin{itemize}
\tightlist
\item
  Fewer layers, cleaner separation of concerns
\item
  User plane: PDCP \ensuremath{\to} RLC \ensuremath{\to} MAC \ensuremath{\to} PHY
\item
  Control plane: NAS (between UE and MME) + RRC (between UE and eNB)
\end{itemize}

\CNsubsection{}{Shared Channel Access}

\begin{itemize}
\tightlist
\item
  No dedicated channels — all resources shared dynamically
\item
  Scheduler in eNodeB makes 1~ms TTI decisions
\item
  Enables statistical multiplexing gains
\end{itemize}

\CNsubsection{}{Self-Organizing Network (SON) Features}

\begin{itemize}
\tightlist
\item
  Automatic neighbor relations (ANR)
\item
  Self-configuration of new eNodeBs
\item
  Self-optimization and self-healing capabilities
\end{itemize}

\CNsection{5.3}{LTE Architecture Overview}

LTE architecture consists of two main domains:

\begin{enumerate}
\def\labelenumi{\arabic{enumi}.}
\tightlist
\item
  \textbf{E-UTRAN} (Evolved UTRAN) — the radio access network
\item
  \textbf{EPC} (Evolved Packet Core) — the core network
\end{enumerate}

\CNfigure{m05-lte-architecture}{LTE architecture: UE, E-UTRAN, EPC and external networks}

\CNsubsection{}{E-UTRAN: The Radio Access Network}

\begin{CNinsight}

\textbf{Key insight: E-UTRAN consists of eNodeBs only — there is NO RNC!}

\end{CNinsight}

\begin{itemize}
\tightlist
\item
  The eNodeB handles ALL radio-related functions previously split between NodeB and RNC
\item
  eNodeBs connect directly to the EPC via the S1 interface
\item
  eNodeBs connect to each other via the X2 interface for handover and load balancing
\item
  This is the ``flat architecture'' — one node type in the RAN
\end{itemize}

\Needspace{19\baselineskip}

\CNsubsection{}{EPC: The Evolved Packet Core}

The EPC is a pure packet core with separated control and user planes:

{\def\LTcaptype{none} % do not increment counter
\Needspace{13\baselineskip}
\begin{longtable}[c]{@{}cc@{}}
\toprule\noalign{}
\rowcolor{CNink}\CNth{Element} & \CNth{Primary Role} \\
\CNheadrulesep
\endhead
\bottomrule\noalign{}
\endlastfoot
\textbf{MME} & Control plane — signaling, mobility, security \\
\textbf{S-GW} & User plane anchor within E-UTRAN \\
\textbf{P-GW} & User plane anchor to external networks \\
\textbf{HSS} & Subscriber database and authentication \\
\textbf{PCRF} & Policy and charging rules \\
\end{longtable}
}

\Needspace{14\baselineskip}

\CNsection{5.4}{Network Elements — Detailed}

\CNchained\CNsubsection{5.4.1}{UE (User Equipment)}

\textbf{Purpose:} The mobile device that communicates over the LTE air interface.

\textbf{Components:}

\begin{itemize}
\tightlist
\item
  \textbf{ME (Mobile Equipment):} The physical device (phone, tablet, modem)
\item
  \textbf{USIM (Universal Subscriber Identity Module):} Stores subscriber identity (IMSI), security keys (K), and authentication algorithms
\end{itemize}

\Needspace{23\baselineskip}

\textbf{UE Categories (3GPP Release 8–12):}

{\def\LTcaptype{none} % do not increment counter
\Needspace{20\baselineskip}
\begin{longtable}[c]{@{}ccccc@{}}
\toprule\noalign{}
\rowcolor{CNink}\CNth{Category} & \CNth{Max DL (Mbps)} & \CNth{Max UL (Mbps)} & \CNth{DL MIMO} & \CNth{Notes} \\
\CNheadrulesep
\endhead
\bottomrule\noalign{}
\endlastfoot
Cat 1 & 10 & 5 & None & IoT, basic \\
Cat 2 & 51 & 25 & 2\ensuremath{\times}2 & Entry smartphones \\
Cat 3 & 100 & 50 & 2\ensuremath{\times}2 & Mainstream \\
Cat 4 & 150 & 50 & 2\ensuremath{\times}2 & High-end \\
Cat 5 & 300 & 75 & 4\ensuremath{\times}4 & Flagship \\
Cat 6 & 300 & 50 & 2\ensuremath{\times}2 & Carrier Aggregation \\
Cat 9 & 450 & 50 & 2\ensuremath{\times}2 & 3CC CA \\
Cat 12 & 600 & 100 & 4\ensuremath{\times}4 & Advanced \\
Cat 16 & 980 & 150 & 4\ensuremath{\times}4 & LTE-Advanced Pro \\
\end{longtable}
}

\textbf{Key UE Capabilities:}

\begin{itemize}
\tightlist
\item
  Carrier Aggregation (CA): combine multiple component carriers
\item
  MIMO support (2\ensuremath{\times}2 or 4\ensuremath{\times}4 antenna configurations)
\item
  VoLTE support (IMS registration, SIP signaling)
\item
  Dual connectivity (connect to two eNodeBs simultaneously — Release 12+)
\end{itemize}

\CNsubsection{5.4.2}{eNodeB (Evolved NodeB)}

\textbf{Purpose:} The single RAN node that handles ALL radio access functions — replaces both NodeB AND RNC from 3G.

\textbf{Problem it solves:} Eliminates the RNC bottleneck, reduces latency by one hop, enables faster scheduling decisions.

\textbf{Functions absorbed from 3G RNC:}

\begin{itemize}
\tightlist
\item
  Radio Resource Management (RRM) — scheduling, admission control
\item
  Radio Bearer Control — setup, maintenance, release of bearers
\item
  Connection Mobility Control — handover decisions
\item
  Dynamic Resource Allocation (Scheduler) — 1~ms TTI granularity
\item
  Inter-cell Interference Coordination (ICIC)
\item
  Header compression (ROHC via PDCP)
\item
  Encryption and integrity protection (PDCP layer)
\item
  Segmentation and reassembly (RLC layer)
\end{itemize}

\textbf{Functions unique to eNodeB:}

\begin{itemize}
\tightlist
\item
  S1 interface termination (to MME and S-GW)
\item
  X2 interface termination (to neighboring eNodeBs)
\item
  Paging message distribution
\item
  Broadcast information (SIBs)
\item
  Measurement configuration and reporting
\end{itemize}

\textbf{Scheduler — The Brain of eNodeB:}

\begin{itemize}
\tightlist
\item
  Makes resource allocation decisions every 1~ms (TTI)
\item
  Assigns Resource Blocks (RBs) to UEs based on:

  \begin{itemize}
  \tightlist
  \item
    Channel quality (CQI reports from UEs)
  \item
    QoS requirements (GBR vs non-GBR bearers)
  \item
    Fairness algorithms (Proportional Fair, Round Robin, Max Throughput)
  \item
    Buffer status and priority
  \end{itemize}
\end{itemize}

\CNsubsection{5.4.3}{MME (Mobility Management Entity)}

\textbf{Purpose:} The control plane anchor of the EPC — handles all signaling and control functions, touches NO user data.

\Needspace{26\baselineskip}

\textbf{Problem it solves:} Separates control from user plane, enabling independent scaling. Centralizes mobility and security management.

\textbf{Key Functions:}

{\def\LTcaptype{none} % do not increment counter
\Needspace{21\baselineskip}
\begin{longtable}[c]{@{}
  >{\raggedright\arraybackslash\hspace{0pt}}m{(\linewidth - 2\tabcolsep) * \real{0.4348}}
  >{\raggedright\arraybackslash\hspace{0pt}}m{(\linewidth - 2\tabcolsep) * \real{0.5652}}@{}}
\toprule\noalign{}
\rowcolor{CNink}\CNth{Function} & \CNth{Description} \\
\CNheadrulesep
\endhead
\bottomrule\noalign{}
\endlastfoot
NAS Signaling & Manages Non-Access Stratum protocols (EMM, ESM) between UE and MME \\
Authentication & Generates authentication vectors with HSS (EPS-AKA) \\
Security & NAS encryption/\allowbreak{}integrity; derives keys for AS security \\
Mobility Management & Tracks UE location (Tracking Areas), handles TAU \\
Bearer Management & Establishes, modifies, releases EPS bearers \\
Paging & Initiates paging when downlink data arrives for idle UE \\
S-GW Selection & Selects appropriate Serving Gateway for UE \\
P-GW Selection & Selects PDN Gateway (during initial attach or new PDN connection) \\
Handover Control & Inter-eNB and inter-MME handover signaling \\
Idle Mode & Manages UE state transitions (connected \ensuremath{\leftrightarrow} idle) \\
\end{longtable}
}

\textbf{NAS Procedures:}

\begin{itemize}
\tightlist
\item
  \textbf{Attach:} UE registers with the network (identity, authentication, bearer setup)
\item
  \textbf{Detach:} UE deregisters (voluntary or network-initiated)
\item
  \textbf{TAU (Tracking Area Update):} UE reports new location when crossing TA boundary
\item
  \textbf{Service Request:} Idle UE requests to become active (for data transfer)
\end{itemize}

\CNsubsection{5.4.4}{S-GW (Serving Gateway)}

\textbf{Purpose:} The user plane anchor within the E-UTRAN — all user data passes through S-GW.

\textbf{Problem it solves:} Provides a stable user plane anchor point during intra-LTE mobility (handover between eNodeBs). Only one tunnel switch needed at S-GW during handover instead of rerouting through external networks.

\textbf{Key Functions:}

\begin{itemize}
\tightlist
\item
  \textbf{Local mobility anchor:} During inter-eNB handover, only the S1-U tunnel is switched at S-GW
\item
  \textbf{Packet routing and forwarding:} Routes packets between eNB and P-GW
\item
  \textbf{Data buffering during handover:} Buffers DL packets while handover completes
\item
  \textbf{Downlink data notification:} Triggers paging via MME when DL data arrives for idle UE
\item
  \textbf{Lawful intercept:} Interface for legal interception of user traffic
\item
  \textbf{Charging data collection:} Usage data for offline charging (CDRs)
\item
  \textbf{Transport level marking:} DSCP marking for QoS in transport network
\item
  \textbf{Idle mode buffering:} Buffers first DL packet and notifies MME for paging
\end{itemize}

\textbf{Deployment Note:} In many networks, S-GW and P-GW are combined into a single node (SAE-GW) for simplicity.

\CNsubsection{5.4.5}{P-GW (PDN Gateway)}

\textbf{Purpose:} The point of connection to external packet data networks (Internet, IMS, enterprise VPNs). Allocates IP addresses and enforces policy.

\textbf{Problem it solves:} Provides a single, stable IP anchor point for the UE regardless of mobility. The UE keeps the same IP address even when moving between S-GWs.

\textbf{Key Functions:}

\begin{itemize}
\tightlist
\item
  \textbf{IP address allocation:} Assigns IPv4 and/\allowbreak{}or IPv6 addresses to UEs (via DHCPv4, SLAAC, or static)
\item
  \textbf{Policy enforcement:} Applies QoS rules received from PCRF (via Gx interface)
\item
  \textbf{Charging:} Online charging (Gy to OCS) and offline charging (Gz to OFCS)
\item
  \textbf{Packet filtering:} Deep packet inspection, TFT (Traffic Flow Template) enforcement
\item
  \textbf{Per-SDF (Service Data Flow) treatment:} Apply different QoS/\allowbreak{}charging per traffic flow
\item
  \textbf{Inter-operator mobility anchor:} When roaming, P-GW remains in home network (home-routed) or visited network (local breakout)
\item
  \textbf{NAT (if needed):} Network Address Translation for IPv4
\item
  \textbf{GTP/\allowbreak{}PMIP anchor:} Terminates S5/\allowbreak{}S8 tunnel from S-GW
\end{itemize}

\textbf{APN (Access Point Name):}

\begin{itemize}
\tightlist
\item
  Each PDN connection is identified by an APN
\item
  P-GW selection based on APN (e.g., ``internet'', ``ims'', ``enterprise.vpn'')
\item
  Multiple simultaneous PDN connections possible (e.g., Internet + IMS for VoLTE)
\end{itemize}

\CNsubsection{5.4.6}{HSS (Home Subscriber Server)}

\textbf{Purpose:} The central database containing all subscriber information and security credentials. Evolution of 3G HLR/\allowbreak{}AuC.

\textbf{Problem it solves:} Centralizes subscriber management, provides authentication vectors, and stores subscription data (allowed APNs, QoS profiles, roaming permissions).

\textbf{Key Functions:}

\begin{itemize}
\tightlist
\item
  \textbf{Subscriber data storage:}

  \begin{itemize}
  \tightlist
  \item
    IMSI (permanent identity)
  \item
    MSISDN (phone number)
  \item
    Subscribed APNs and associated QoS profiles
  \item
    Roaming restrictions
  \item
    Subscribed Tracking Areas
  \end{itemize}
\item
  \textbf{Authentication vector generation:}

  \begin{itemize}
  \tightlist
  \item
    Stores permanent key K (shared with USIM)
  \item
    Generates authentication vectors (RAND, AUTN, XRES, \ensuremath{K_{\mathrm{ASME}}}) using EPS-AKA
  \item
    Provides vectors to MME for UE authentication
  \end{itemize}
\item
  \textbf{Location management:}

  \begin{itemize}
  \tightlist
  \item
    Stores current serving MME for each subscriber
  \item
    Used for terminating calls/\allowbreak{}sessions (where to page the UE)
  \end{itemize}
\item
  \textbf{Authorization:}

  \begin{itemize}
  \tightlist
  \item
    Determines which services a subscriber can access
  \item
    Enforces subscription limits (max bitrate, allowed APNs)
  \end{itemize}
\end{itemize}

\textbf{Interfaces:}

\begin{itemize}
\tightlist
\item
  S6a to MME (Diameter protocol)
\item
  Cx/\allowbreak{}Dx to IMS (for VoLTE subscriber data)
\item
  Sh to Application Servers
\end{itemize}

\CNsubsection{5.4.7}{PCRF (Policy and Charging Rules Function)}

\textbf{Purpose:} The policy brain of the network — makes real-time decisions about QoS and charging rules applied to each data flow.

\textbf{Problem it solves:} Enables dynamic, per-flow QoS (e.g., guarantee bandwidth for a VoLTE call while best-effort for web browsing). Enables flexible charging models (zero-rating, tiered plans, sponsored data).

\textbf{Key Functions:}

\begin{itemize}
\tightlist
\item
  \textbf{Policy decision:} Determines QoS parameters (QCI, ARP, MBR, GBR) for each service data flow
\item
  \textbf{Charging rules:} Specifies how each flow is charged (rating group, metering method)
\item
  \textbf{PCC (Policy and Charging Control) rules:} Combined policy + charging instructions sent to P-GW
\item
  \textbf{Dynamic policy updates:} Can modify rules in real-time (e.g., when user starts a video call)
\item
  \textbf{Subscription-based policy:} Considers user’s subscription tier
\item
  \textbf{Network-condition-based policy:} Can throttle based on congestion
\end{itemize}

\textbf{Example PCC Rule:}

\tcbset{CNcodemode/.style={unbreakable}}\def\CNcodelang{}

\begin{CNplaincode}
\begin{Verbatim}
Rule: "VoLTE_Voice"
  - Service Data Flow: SIP/RTP traffic to IMS
  - QCI: 1 (Conversational Voice)
  - GBR: 40 kbps UL / 40 kbps DL
  - ARP: Priority 1, pre-emption capable
  - Charging: Included in voice plan (no data deduction)
\end{Verbatim}
\end{CNplaincode}

\textbf{Interfaces:}

\begin{itemize}
\tightlist
\item
  Gx to P-GW (installs/\allowbreak{}modifies PCC rules)
\item
  Rx from IMS/\allowbreak{}AF (receives session info, triggers policy)
\item
  Sp to SPR (Subscription Profile Repository)
\end{itemize}

\Needspace{14\baselineskip}

\CNkeepfig{m05-interface-map}{8}

\CNsection{5.5}{LTE Interfaces — Reference Points}

\CNchained\CNsubsection{}{Interface Map}

\CNfigure{m05-interface-map}{LTE interface map with the protocol of each reference point}

\Needspace{15\baselineskip}

\CNsubsection{5.5.1}{S1-MME Interface}

{\def\LTcaptype{none} % do not increment counter
\Needspace{11\baselineskip}
\begin{longtable}[c]{@{}cc@{}}
\toprule\noalign{}
\rowcolor{CNink}\CNth{Property} & \CNth{Value} \\
\CNheadrulesep
\endhead
\bottomrule\noalign{}
\endlastfoot
\textbf{Endpoints} & eNodeB \ensuremath{\leftrightarrow} MME \\
\textbf{Plane} & Control Plane \\
\textbf{Protocol Stack} & S1AP /\allowbreak{} SCTP /\allowbreak{} IP \\
\textbf{Purpose} & All control signaling between RAN and core \\
\end{longtable}
}

\textbf{What flows over S1-MME:}

\begin{itemize}
\tightlist
\item
  Initial UE messages (Attach, TAU, Service Request)
\item
  UE context setup/\allowbreak{}release
\item
  Handover preparation and execution (inter-eNB via MME)
\item
  Paging requests from MME to eNB
\item
  Bearer setup/\allowbreak{}modification/\allowbreak{}release commands
\item
  NAS message transport (MME\ensuremath{\leftrightarrow}UE via eNB as relay)
\item
  Reset and error indication procedures
\end{itemize}

\textbf{Why SCTP (not TCP):}

\begin{itemize}
\tightlist
\item
  Multi-homing: survives link failures without session drop
\item
  Multi-streaming: avoids head-of-line blocking
\item
  Message-oriented (not byte-stream): natural fit for signaling
\item
  Built-in heartbeat for connection monitoring
\end{itemize}

\Needspace{15\baselineskip}

\CNsubsection{5.5.2}{S1-U Interface}

{\def\LTcaptype{none} % do not increment counter
\Needspace{11\baselineskip}
\begin{longtable}[c]{@{}cc@{}}
\toprule\noalign{}
\rowcolor{CNink}\CNth{Property} & \CNth{Value} \\
\CNheadrulesep
\endhead
\bottomrule\noalign{}
\endlastfoot
\textbf{Endpoints} & eNodeB \ensuremath{\leftrightarrow} S-GW \\
\textbf{Plane} & User Plane \\
\textbf{Protocol Stack} & GTP-U /\allowbreak{} UDP /\allowbreak{} IP \\
\textbf{Purpose} & Carries user IP packets in GTP tunnels \\
\end{longtable}
}

\textbf{What flows over S1-U:}

\begin{itemize}
\tightlist
\item
  All user data (encapsulated in GTP-U tunnels)
\item
  One GTP tunnel per EPS bearer per UE
\item
  End marker packets during handover (signals last packet from old path)
\end{itemize}

\textbf{GTP-U Tunnel Identification:}

\begin{itemize}
\tightlist
\item
  Each tunnel identified by TEID (Tunnel Endpoint Identifier) + IP address
\item
  TEIDs allocated independently at each end
\item
  Enables multiplexing many UE bearers over single transport connection
\end{itemize}

\textbf{Why GTP-U (not plain IP):}

\begin{itemize}
\tightlist
\item
  Tunneling allows mobility transparency (IP address doesn’t change)
\item
  Per-bearer QoS enforcement
\item
  Simple encapsulation: GTP-U header (8+ bytes) + inner IP packet
\end{itemize}

\Needspace{15\baselineskip}

\CNsubsection{5.5.3}{X2 Interface}

{\def\LTcaptype{none} % do not increment counter
\Needspace{11\baselineskip}
\begin{longtable}[c]{@{}cc@{}}
\toprule\noalign{}
\rowcolor{CNink}\CNth{Property} & \CNth{Value} \\
\CNheadrulesep
\endhead
\bottomrule\noalign{}
\endlastfoot
\textbf{Endpoints} & eNodeB \ensuremath{\leftrightarrow} eNodeB \\
\textbf{Plane} & Control + User Plane \\
\textbf{Protocol Stack} & X2AP/\allowbreak{}SCTP (control) + GTP-U/\allowbreak{}UDP (user) \\
\textbf{Purpose} & Direct inter-eNB handover and load management \\
\end{longtable}
}

\textbf{Control Plane (X2AP) Functions:}

\begin{itemize}
\tightlist
\item
  Handover preparation (request, acknowledge, cancel)
\item
  SN (Sequence Number) status transfer
\item
  Load indication (exchange load information between eNBs)
\item
  Resource status reporting
\item
  Inter-cell interference coordination (ICIC) information exchange
\end{itemize}

\textbf{User Plane (GTP-U) Functions:}

\begin{itemize}
\tightlist
\item
  Forward DL data from source eNB to target eNB during handover
\item
  Prevents data loss during handover gap
\item
  Temporary tunnel — torn down after handover completes
\end{itemize}

\textbf{X2 Handover Advantage:}

\begin{itemize}
\tightlist
\item
  Direct eNB-to-eNB signaling — no need to involve MME for preparation
\item
  Only MME notified after handover completes (path switch)
\item
  Faster than S1-based handover (used when X2 is available)
\end{itemize}

\Needspace{15\baselineskip}

\CNsubsection{5.5.4}{S5/\allowbreak{}S8 Interface}

{\def\LTcaptype{none} % do not increment counter
\Needspace{11\baselineskip}
\begin{longtable}[c]{@{}cc@{}}
\toprule\noalign{}
\rowcolor{CNink}\CNth{Property} & \CNth{Value} \\
\CNheadrulesep
\endhead
\bottomrule\noalign{}
\endlastfoot
\textbf{Endpoints} & S-GW \ensuremath{\leftrightarrow} P-GW \\
\textbf{Plane} & User + Control Plane \\
\textbf{Protocol Stack} & GTP-C + GTP-U (or PMIP alternative) \\
\textbf{Purpose} & Connects serving gateway to PDN gateway \\
\end{longtable}
}

\textbf{S5 vs S8:}

\begin{itemize}
\tightlist
\item
  \textbf{S5:} When S-GW and P-GW are in the \textbf{same} PLMN (non-roaming or home-routed)
\item
  \textbf{S8:} When S-GW and P-GW are in \textbf{different} PLMNs (roaming scenario)
\item
  Technically identical protocol, different name indicates inter-PLMN boundary
\end{itemize}

\textbf{Functions:}

\begin{itemize}
\tightlist
\item
  Tunnel management (create, update, delete sessions)
\item
  Per-bearer GTP tunnels for user plane
\item
  Triggered by MME (via S11) during attach or bearer modification
\item
  GTP-C handles session/\allowbreak{}bearer management signaling
\item
  GTP-U carries encapsulated user data
\end{itemize}

\textbf{Protocol Options:}

\begin{itemize}
\tightlist
\item
  \textbf{GTP-based S5/\allowbreak{}S8:} Full GTP-C signaling + GTP-U tunnels (most common)
\item
  \textbf{PMIP-based S5/\allowbreak{}S8:} Proxy Mobile IPv6 signaling (rare in practice)
\end{itemize}

\Needspace{15\baselineskip}

\CNsubsection{5.5.5}{S6a Interface}

{\def\LTcaptype{none} % do not increment counter
\Needspace{11\baselineskip}
\begin{longtable}[c]{@{}cc@{}}
\toprule\noalign{}
\rowcolor{CNink}\CNth{Property} & \CNth{Value} \\
\CNheadrulesep
\endhead
\bottomrule\noalign{}
\endlastfoot
\textbf{Endpoints} & MME \ensuremath{\leftrightarrow} HSS \\
\textbf{Plane} & Control Plane \\
\textbf{Protocol Stack} & Diameter /\allowbreak{} SCTP (or TCP) /\allowbreak{} IP \\
\textbf{Purpose} & Subscriber authentication and data retrieval \\
\end{longtable}
}

\Needspace{21\baselineskip}

\textbf{Diameter Commands on S6a:}

{\def\LTcaptype{none} % do not increment counter
\Needspace{18\baselineskip}
\begin{longtable}[c]{@{}
  >{\raggedright\arraybackslash\hspace{0pt}}m{(\linewidth - 4\tabcolsep) * \real{0.3103}}
  >{\raggedright\arraybackslash\hspace{0pt}}m{(\linewidth - 4\tabcolsep) * \real{0.3793}}
  >{\raggedright\arraybackslash\hspace{0pt}}m{(\linewidth - 4\tabcolsep) * \real{0.3103}}@{}}
\toprule\noalign{}
\rowcolor{CNink}\CNth{Command} & \CNth{Direction} & \CNth{Purpose} \\
\CNheadrulesep
\endhead
\bottomrule\noalign{}
\endlastfoot
Authentication-Information-Request (AIR) & MME\ensuremath{\to}HSS & Request auth vectors \\
Authentication-Information-Answer (AIA) & HSS\ensuremath{\to}MME & Return auth vectors (RAND, AUTN, XRES, \ensuremath{K_{\mathrm{ASME}}}) \\
Update-Location-Request (ULR) & MME\ensuremath{\to}HSS & Register MME as serving node \\
Update-Location-Answer (ULA) & HSS\ensuremath{\to}MME & Return subscription data \\
Cancel-Location-Request (CLR) & HSS\ensuremath{\to}MME & Deregister UE from old MME \\
Insert-Subscriber-Data-Request (IDR) & HSS\ensuremath{\to}MME & Push updated subscription data \\
Delete-Subscriber-Data-Request (DSR) & HSS\ensuremath{\to}MME & Remove subscription data \\
Purge-UE-Request (PUR) & MME\ensuremath{\to}HSS & Inform HSS that UE is unreachable \\
\end{longtable}
}

\textbf{Why Diameter (not SS7/\allowbreak{}MAP):}

\begin{itemize}
\tightlist
\item
  IP-native (aligns with all-IP architecture)
\item
  Extensible via AVPs (Attribute-Value Pairs)
\item
  Better security (TLS/\allowbreak{}IPsec support)
\item
  Larger message sizes (supports complex subscription data)
\end{itemize}

\Needspace{15\baselineskip}

\CNsubsection{5.5.6}{S11 Interface}

{\def\LTcaptype{none} % do not increment counter
\Needspace{11\baselineskip}
\begin{longtable}[c]{@{}cc@{}}
\toprule\noalign{}
\rowcolor{CNink}\CNth{Property} & \CNth{Value} \\
\CNheadrulesep
\endhead
\bottomrule\noalign{}
\endlastfoot
\textbf{Endpoints} & MME \ensuremath{\leftrightarrow} S-GW \\
\textbf{Plane} & Control Plane \\
\textbf{Protocol Stack} & GTP-C v2 /\allowbreak{} UDP /\allowbreak{} IP \\
\textbf{Purpose} & EPS bearer/\allowbreak{}session management between MME and S-GW \\
\end{longtable}
}

\textbf{Key Messages:}

\begin{itemize}
\tightlist
\item
  \textbf{Create Session Request/\allowbreak{}Response:} During attach — establishes default bearer
\item
  \textbf{Modify Bearer Request/\allowbreak{}Response:} During handover — updates S1-U tunnel info
\item
  \textbf{Delete Session Request/\allowbreak{}Response:} During detach — tears down bearers
\item
  \textbf{Create Bearer Request/\allowbreak{}Response:} MME-initiated dedicated bearer setup
\item
  \textbf{Release Access Bearers:} When UE goes idle (release S1-U but keep S5/\allowbreak{}S8)
\item
  \textbf{Downlink Data Notification:} S-GW tells MME that data arrived for idle UE \ensuremath{\to} triggers paging
\end{itemize}

\textbf{Why GTP-C v2:}

\begin{itemize}
\tightlist
\item
  Already proven in GPRS/\allowbreak{}3G (GTP-C v1)
\item
  Efficient binary encoding
\item
  Built-in TEID-based message routing
\item
  Supports piggybacking (multiple IEs in one message)
\end{itemize}

\Needspace{15\baselineskip}

\CNsubsection{5.5.7}{SGi Interface}

{\def\LTcaptype{none} % do not increment counter
\Needspace{11\baselineskip}
\begin{longtable}[c]{@{}cc@{}}
\toprule\noalign{}
\rowcolor{CNink}\CNth{Property} & \CNth{Value} \\
\CNheadrulesep
\endhead
\bottomrule\noalign{}
\endlastfoot
\textbf{Endpoints} & P-GW \ensuremath{\leftrightarrow} External Networks (Internet, IMS, Enterprise) \\
\textbf{Plane} & User Plane \\
\textbf{Protocol Stack} & IP (standard Internet protocols) \\
\textbf{Purpose} & Connects LTE network to the outside world \\
\end{longtable}
}

\textbf{Characteristics:}

\begin{itemize}
\tightlist
\item
  Equivalent to the Gi interface in 2G/\allowbreak{}3G GPRS
\item
  Standard IP interface — no GTP or mobile-specific protocols
\item
  P-GW performs NAT here (if needed) or routes native IPv6
\item
  This is where the UE’s IP address is ``visible'' to the internet
\item
  Firewall, DPI, and value-added services typically deployed here
\end{itemize}

\textbf{Connected Networks:}

\begin{itemize}
\tightlist
\item
  Public Internet
\item
  IMS core (for VoLTE, VoWiFi)
\item
  Enterprise VPNs (via APN-based routing)
\item
  Content Delivery Networks (CDNs)
\item
  Walled garden services
\end{itemize}

\Needspace{15\baselineskip}

\CNsubsection{5.5.8}{Gx Interface}

{\def\LTcaptype{none} % do not increment counter
\Needspace{11\baselineskip}
\begin{longtable}[c]{@{}
  >{\raggedright\arraybackslash\hspace{0pt}}m{(\linewidth - 2\tabcolsep) * \real{0.5882}}
  >{\raggedright\arraybackslash\hspace{0pt}}m{(\linewidth - 2\tabcolsep) * \real{0.4118}}@{}}
\toprule\noalign{}
\rowcolor{CNink}\CNth{Property} & \CNth{Value} \\
\CNheadrulesep
\endhead
\bottomrule\noalign{}
\endlastfoot
\textbf{Endpoints} & PCRF \ensuremath{\leftrightarrow} P-GW \\
\textbf{Plane} & Control Plane \\
\textbf{Protocol Stack} & Diameter /\allowbreak{} SCTP or TCP /\allowbreak{} IP \\
\textbf{Purpose} & Install, modify, remove PCC (Policy and Charging Control) rules \\
\end{longtable}
}

\textbf{How it works:}

\begin{enumerate}
\def\labelenumi{\arabic{enumi}.}
\tightlist
\item
  P-GW contacts PCRF when a new IP-CAN session is established (UE attaches)
\item
  PCRF evaluates subscription data + network policy + application input
\item
  PCRF sends PCC rules to P-GW specifying per-flow QoS and charging
\item
  PCRF can push updated rules at any time (e.g., when VoLTE call starts)
\end{enumerate}

\textbf{Diameter Commands:}

\begin{itemize}
\tightlist
\item
  \textbf{CC-Request/\allowbreak{}Answer (CCR/\allowbreak{}CCA):} P-GW reports events, PCRF responds with rules
\item
  \textbf{Re-Auth-Request/\allowbreak{}Answer (RAR/\allowbreak{}RAA):} PCRF pushes new rules to P-GW
\end{itemize}

\textbf{PCC Rule Contents:}

\begin{itemize}
\tightlist
\item
  Service Data Flow (SDF) filter (5-tuple: src/\allowbreak{}dst IP, ports, protocol)
\item
  QoS parameters (QCI, ARP, MBR, GBR)
\item
  Charging key (rating group)
\item
  Metering method (volume, time, event)
\item
  Gate status (open/\allowbreak{}close a flow)
\end{itemize}

\Needspace{15\baselineskip}

\CNsubsection{5.5.9}{Gy Interface}

{\def\LTcaptype{none} % do not increment counter
\Needspace{11\baselineskip}
\begin{longtable}[c]{@{}cc@{}}
\toprule\noalign{}
\rowcolor{CNink}\CNth{Property} & \CNth{Value} \\
\CNheadrulesep
\endhead
\bottomrule\noalign{}
\endlastfoot
\textbf{Endpoints} & P-GW \ensuremath{\leftrightarrow} OCS (Online Charging System) \\
\textbf{Plane} & Control Plane (charging signaling) \\
\textbf{Protocol Stack} & Diameter /\allowbreak{} SCTP or TCP /\allowbreak{} IP \\
\textbf{Purpose} & Real-time credit control for prepaid/\allowbreak{}online charging \\
\end{longtable}
}

\textbf{How it works:}

\begin{enumerate}
\def\labelenumi{\arabic{enumi}.}
\tightlist
\item
  When a prepaid user starts a data session, P-GW requests quota from OCS
\item
  OCS grants units (bytes, seconds, or events)
\item
  P-GW meters usage against granted quota
\item
  When quota runs low, P-GW requests more (or session is terminated)
\item
  Enables real-time balance deduction and spend limits
\end{enumerate}

\textbf{Diameter Commands:}

\begin{itemize}
\tightlist
\item
  \textbf{Credit-Control-Request (CCR):} P-GW requests/\allowbreak{}reports usage

  \begin{itemize}
  \tightlist
  \item
    INITIAL: Session start, request first quota
  \item
    UPDATE: Quota running low, request more
  \item
    TERMINATION: Session end, report final usage
  \end{itemize}
\item
  \textbf{Credit-Control-Answer (CCA):} OCS grants quota or denies
\end{itemize}

\textbf{Use Cases:}

\begin{itemize}
\tightlist
\item
  Prepaid data plans with real-time balance
\item
  Fair usage policies (throttle after threshold)
\item
  Sponsored data (third party pays)
\item
  Zero-rating specific applications
\end{itemize}

\Needspace{14\baselineskip}

\CNsection{5.6}{IMS Overview (IP Multimedia Subsystem)}

\CNchained\CNsubsection{}{Why VoLTE Needs IMS}

LTE is a \textbf{pure packet network} — there is no circuit-switched domain for voice. Without IMS:

\begin{itemize}
\tightlist
\item
  Voice calls would fall back to 2G/\allowbreak{}3G (CSFB — Circuit-Switched Fallback)
\item
  CSFB adds 2–4 seconds call setup delay and drops LTE data temporarily
\item
  No way to provide rich communication services (video calls, messaging)
\end{itemize}

\textbf{IMS provides the session control layer that makes VoLTE possible.}

\CNkeepfig{m05-ims-arch}{3}

\CNsubsection{}{IMS Basic Architecture}

\CNfigure{m05-ims-arch}{IMS signalling path from the UE to the application servers}

\textbf{Key IMS Elements:}\CNbr{}\textbar{} Element \textbar{} Role \textbar{}\CNbr{}\textbar{}———\textbar{}——\textbar{}\CNbr{}\textbar{} \textbf{P-CSCF} (Proxy-CSCF) \textbar{} First contact point; SIP proxy in visited network; provides security (IPsec) \textbar{}\CNbr{}\textbar{} \textbf{I-CSCF} (Interrogating-CSCF) \textbar{} Entry point in home network; queries HSS for S-CSCF assignment \textbar{}\CNbr{}\textbar{} \textbf{S-CSCF} (Serving-CSCF) \textbar{} Core SIP registrar and session controller; applies service logic \textbar{}\CNbr{}\textbar{} \textbf{Application Servers} \textbar{} VoLTE TAS (Telephony AS), conference, messaging services \textbar{}\CNbr{}\textbar{} \textbf{MGCF/\allowbreak{}MGW} \textbar{} Media Gateway — interworks with PSTN for legacy calls \textbar{}

\CNsubsection{}{VoLTE Call Flow (Simplified)}

\begin{enumerate}
\def\labelenumi{\arabic{enumi}.}
\tightlist
\item
  UE registers with IMS via SIP REGISTER (through P-CSCF \ensuremath{\to} S-CSCF)
\item
  UE initiates call via SIP INVITE
\item
  IMS signals PCRF (via Rx interface) to request dedicated bearer
\item
  PCRF instructs P-GW (via Gx) to create GBR bearer (QCI=1 for voice)
\item
  MME/\allowbreak{}eNB establish dedicated bearer with guaranteed bitrate
\item
  RTP media flows over the dedicated bearer
\item
  Result: HD voice with guaranteed QoS, \textasciitilde{}100~ms end-to-end
\end{enumerate}

\CNsubsection{}{IMS Integration with LTE}

\begin{itemize}
\tightlist
\item
  P-CSCF discovered by UE during PDN connection to IMS APN
\item
  UE maintains two PDN connections: Internet APN + IMS APN
\item
  Rx interface: IMS \ensuremath{\to} PCRF (triggers dedicated bearer for voice/\allowbreak{}video)
\item
  VoLTE requires: UE support + IMS core + PCRF + dedicated bearers
\end{itemize}

\Needspace{14\baselineskip}

\CNkeepfig{m05-cups}{10}

\CNsection{5.7}{Control Plane vs User Plane Separation (CUPS)}

\CNchained\CNsubsection{}{The Concept}

CUPS (3GPP Release 14, TS 23.214) separates the control and user plane functions of S-GW and P-GW into independent entities that can be scaled and deployed independently.

\CNfigure{m05-cups}{Combined gateways versus CUPS (control/\allowbreak{}user plane separation)}

\Needspace{17\baselineskip}

\CNsubsection{}{Why CUPS Matters}

{\def\LTcaptype{none} % do not increment counter
\Needspace{13\baselineskip}
\begin{longtable}[c]{@{}
  >{\raggedright\arraybackslash\hspace{0pt}}m{(\linewidth - 4\tabcolsep) * \real{0.2500}}
  >{\raggedright\arraybackslash\hspace{0pt}}m{(\linewidth - 4\tabcolsep) * \real{0.4062}}
  >{\raggedright\arraybackslash\hspace{0pt}}m{(\linewidth - 4\tabcolsep) * \real{0.3438}}@{}}
\toprule\noalign{}
\rowcolor{CNink}\CNth{Aspect} & \CNth{Without CUPS} & \CNth{With CUPS} \\
\CNheadrulesep
\endhead
\bottomrule\noalign{}
\endlastfoot
Scaling & Scale entire gateway (wasteful) & Scale UP and CP independently \\
Deployment & GW in central DC only & UP at edge, CP centralized \\
Latency & All traffic through central GW & UP at edge = lower latency \\
Cost & Expensive combined HW & Cheaper, commodity UP hardware \\
Evolution & Monolithic & Direct path to 5G UPF architecture \\
\end{longtable}
}

\CNsubsection{}{CUPS Interface: Sx (PFCP)}

\begin{itemize}
\tightlist
\item
  \textbf{Protocol:} PFCP (Packet Forwarding Control Protocol)
\item
  \textbf{Sxa:} SGW-C \ensuremath{\leftrightarrow} SGW-U
\item
  \textbf{Sxb:} PGW-C \ensuremath{\leftrightarrow} PGW-U
\item
  \textbf{Sxc:} TDF-C \ensuremath{\leftrightarrow} TDF-U
\item
  CP installs forwarding rules (PDRs, FARs, QERs, URRs) into UP
\item
  UP executes packet forwarding based on installed rules
\end{itemize}

\CNsubsection{}{CUPS as Bridge to 5G}

CUPS is the architectural precursor to 5G’s fully separated architecture:

\begin{itemize}
\tightlist
\item
  SGW-C + PGW-C \ensuremath{\to} 5G SMF (Session Management Function)
\item
  SGW-U + PGW-U \ensuremath{\to} 5G UPF (User Plane Function)
\item
  PFCP protocol reused in 5G (N4 interface = Sx evolution)
\end{itemize}

\Needspace{28\baselineskip}

\CNsection{5.8}{Key Differences from 3G}

{\def\LTcaptype{none} % do not increment counter
\Needspace{22\baselineskip}
\begin{longtable}[c]{@{}
  >{\raggedright\arraybackslash\hspace{0pt}}m{(\linewidth - 4\tabcolsep) * \real{0.2759}}
  >{\raggedright\arraybackslash\hspace{0pt}}m{(\linewidth - 4\tabcolsep) * \real{0.4483}}
  >{\raggedright\arraybackslash\hspace{0pt}}m{(\linewidth - 4\tabcolsep) * \real{0.2759}}@{}}
\toprule\noalign{}
\rowcolor{CNink}\CNth{Aspect} & \CNth{3G UMTS/\allowbreak{}HSPA} & \CNth{4G LTE} \\
\CNheadrulesep
\endhead
\bottomrule\noalign{}
\endlastfoot
\textbf{RAN Architecture} & Hierarchical: NodeB \ensuremath{\to} RNC & Flat: eNodeB only \\
\textbf{RAN Controller} & RNC (centralized) & None — functions in eNodeB \\
\textbf{Core Network} & CS domain + PS domain & EPC (PS only, all-IP) \\
\textbf{Voice} & Circuit-switched (MSC) & VoLTE via IMS (packet) \\
\textbf{Handover} & Soft handover (macro diversity) & Hard handover only \\
\textbf{Air Interface} & WCDMA /\allowbreak{} HSPA & OFDMA (DL) /\allowbreak{} SC-FDMA (UL) \\
\textbf{Channel Type} & Dedicated + shared & All shared (no dedicated channels) \\
\textbf{Scheduling} & 2~ms TTI (HSPA) & 1~ms TTI \\
\textbf{Bearer Concept} & RAB + PDP Context & EPS Bearer (unified) \\
\textbf{Max Bandwidth} & 5~MHz fixed & 1.4–20~MHz flexible \\
\textbf{Peak Rate} & 42~Mbps (HSPA+) & 300+ Mbps (Cat 5+) \\
\textbf{Latency (RTT)} & 50–100~ms & 10–20~ms \\
\textbf{Duplexing} & FDD only (most) & FDD + TDD \\
\textbf{Mobility Anchor} & SGSN (control+user) & MME (control) + S-GW (user) — separated \\
\textbf{Subscriber DB} & HLR (SS7/\allowbreak{}MAP) & HSS (Diameter) \\
\textbf{Policy Control} & Limited (PCRF later addition) & Native PCRF (PCC from day one) \\
\end{longtable}
}

\CNkeepfig{m05-3g-vs-4g}{3}

\CNsubsection{}{Architecture Comparison Diagram}

\CNfigure{m05-3g-vs-4g}{3G UMTS versus 4G LTE architecture}

\CNsubsection{}{Key Architectural Wins}

\begin{enumerate}
\def\labelenumi{\arabic{enumi}.}
\tightlist
\item
  \textbf{Elimination of RNC:}

  \begin{itemize}
  \tightlist
  \item
    Removed single point of failure
  \item
    Reduced latency by one hop
  \item
    Enabled distributed scheduling (faster response)
  \item
    Simplified RAN planning and scaling
  \end{itemize}
\item
  \textbf{All-IP /\allowbreak{} No Circuit Switching:}

  \begin{itemize}
  \tightlist
  \item
    Single protocol stack to manage
  \item
    Reduced OPEX (no MSC, MGW, VLR infrastructure)
  \item
    Statistical multiplexing for voice (VoLTE more efficient)
  \item
    Unified charging and policy for all services
  \end{itemize}
\item
  \textbf{EPS Bearer Model:}

  \begin{itemize}
  \tightlist
  \item
    Unified bearer from UE to P-GW (end-to-end QoS)
  \item
    Replaces fragmented 3G model (RAB + PDP Context + separate QoS at each interface)
  \item
    Default bearer always on (always connected experience)
  \item
    Dedicated bearers for specific QoS needs (VoLTE, video)
  \end{itemize}
\item
  \textbf{Separation of Control and User Planes:}

  \begin{itemize}
  \tightlist
  \item
    MME handles only signaling (scales with number of UEs)
  \item
    S-GW handles only data (scales with traffic volume)
  \item
    Independent evolution and scaling
  \end{itemize}
\end{enumerate}

\Needspace{14\baselineskip}

\CNkeepfig{m05-cp-attach}{8}

\CNsection{5.9}{Control Plane vs User Plane Paths}

\CNchained\CNsubsection{}{Control Plane Path}

\CNfigure{m05-cp-attach}{Control-plane path: the simplified attach procedure}

\CNkeepfig{m05-up-path}{3}

\CNsubsection{}{User Plane Path}

\CNfigure{m05-up-path}{User-plane path from the UE to the Internet}

\CNkeepfig{m05-up-stack}{2}

\textbf{User Plane Protocol Stack:}

\CNfigure{m05-up-stack}{User-plane protocol stack of an LTE packet}

\Needspace{28\baselineskip}

\CNsection{5.10}{Interface Summary Table}

{\def\LTcaptype{none} % do not increment counter
\Needspace{22\baselineskip}
\begin{longtable}[c]{@{}
  >{\raggedright\arraybackslash\hspace{0pt}}m{(\linewidth - 8\tabcolsep) * \real{0.2292}}
  >{\raggedright\arraybackslash\hspace{0pt}}m{(\linewidth - 8\tabcolsep) * \real{0.2292}}
  >{\raggedright\arraybackslash\hspace{0pt}}m{(\linewidth - 8\tabcolsep) * \real{0.1458}}
  >{\raggedright\arraybackslash\hspace{0pt}}m{(\linewidth - 8\tabcolsep) * \real{0.2083}}
  >{\raggedright\arraybackslash\hspace{0pt}}m{(\linewidth - 8\tabcolsep) * \real{0.1875}}@{}}
\toprule\noalign{}
\rowcolor{CNink}\CNth{Interface} & \CNth{Endpoints} & \CNth{Plane} & \CNth{Protocol} & \CNth{Purpose} \\
\CNheadrulesep
\endhead
\bottomrule\noalign{}
\endlastfoot
\textbf{Uu} & UE \ensuremath{\leftrightarrow} eNodeB & Control + User & LTE-Uu (OFDMA/\allowbreak{}SC-FDMA) & Air interface — radio communication \\
\textbf{S1-MME} & eNodeB \ensuremath{\leftrightarrow} MME & Control & S1AP /\allowbreak{} SCTP & RAN-Core control signaling (NAS relay, handover, bearer mgmt) \\
\textbf{S1-U} & eNodeB \ensuremath{\leftrightarrow} S-GW & User & GTP-U /\allowbreak{} UDP /\allowbreak{} IP & User data tunneling between RAN and core \\
\textbf{X2} & eNodeB \ensuremath{\leftrightarrow} eNodeB & Control + User & X2AP/\allowbreak{}SCTP + GTP-U/\allowbreak{}UDP & Inter-eNB handover signaling and data forwarding \\
\textbf{S5/\allowbreak{}S8} & S-GW \ensuremath{\leftrightarrow} P-GW & Control + User & GTP-C + GTP-U (or PMIP) & Session/\allowbreak{}bearer management and user data between gateways \\
\textbf{S6a} & MME \ensuremath{\leftrightarrow} HSS & Control & Diameter /\allowbreak{} SCTP & Authentication vectors and subscriber data retrieval \\
\textbf{S11} & MME \ensuremath{\leftrightarrow} S-GW & Control & GTP-C v2 /\allowbreak{} UDP & Bearer/\allowbreak{}session lifecycle management \\
\textbf{SGi} & P-GW \ensuremath{\leftrightarrow} Internet/\allowbreak{}IMS & User & IP (native) & External PDN connectivity — UE traffic exits here \\
\textbf{Gx} & PCRF \ensuremath{\leftrightarrow} P-GW & Control & Diameter & Policy and charging rule installation/\allowbreak{}modification \\
\textbf{Gy} & P-GW \ensuremath{\leftrightarrow} OCS & Control & Diameter & Online/\allowbreak{}real-time credit control (prepaid charging) \\
\textbf{Rx} & IMS/\allowbreak{}AF \ensuremath{\leftrightarrow} PCRF & Control & Diameter & Application-triggered QoS/\allowbreak{}policy requests \\
\textbf{S10} & MME \ensuremath{\leftrightarrow} MME & Control & GTP-C v2 & Inter-MME handover (UE context transfer) \\
\textbf{S3} & MME \ensuremath{\leftrightarrow} SGSN & Control & GTP-C v2 & 3G/\allowbreak{}4G inter-RAT mobility \\
\textbf{SBc} & MME \ensuremath{\leftrightarrow} CBC & Control & SBc-AP /\allowbreak{} SCTP & Cell broadcast (emergency alerts, CMAS/\allowbreak{}ETWS) \\
\end{longtable}
}

\CNboxsection{Key Takeaways}

\begin{CNbox}[unbreakable]{sum}{Key Takeaways}

\begin{enumerate}
\def\labelenumi{\arabic{enumi}.}
\tightlist
\item
  \textbf{LTE = Flat RAN + All-IP Core} — simplicity drives performance
\item
  \textbf{eNodeB is the only RAN node} — absorbed all RNC functions
\item
  \textbf{EPC separates control (MME) from user plane (S-GW/\allowbreak{}P-GW)} — independent scaling
\item
  \textbf{Every interface has a clear, single purpose} — modular, replaceable
\item
  \textbf{GTP tunnels provide mobility transparency} — IP address never changes regardless of movement
\item
  \textbf{PCRF enables dynamic QoS} — different treatment per application flow
\item
  \textbf{IMS replaces circuit-switched voice} — VoLTE = SIP + dedicated bearer
\item
  \textbf{CUPS is the bridge to 5G} — same separation philosophy, evolved into 5G SBA
\end{enumerate}

\end{CNbox}

\CNsection{}{References}

\begin{itemize}
\tightlist
\item
  3GPP TS 23.401: ``GPRS Enhancements for E-UTRAN Access'' (EPC architecture)
\item
  3GPP TS 36.300: ``E-UTRA and E-UTRAN; Overall Description'' (RAN architecture)
\item
  3GPP TS 23.214: ``Architecture Enhancements for CUPS''
\item
  3GPP TS 29.274: ``GTPv2-C'' (S11, S5/\allowbreak{}S8 control)
\item
  3GPP TS 29.281: ``GTP-U'' (S1-U, S5/\allowbreak{}S8 user)
\item
  3GPP TS 36.413: ``S1AP'' (S1-MME interface)
\item
  3GPP TS 36.423: ``X2AP'' (X2 interface)
\item
  3GPP TS 29.272: ``S6a/\allowbreak{}S6d'' (MME-HSS Diameter interface)
\item
  3GPP TS 29.212: ``Gx'' (PCRF-PCEF Diameter interface)
\end{itemize}

\emph{Module 5 Complete — Next: Module 6: 5G NR Architecture}
